% Transcription — fonds Alexandre Grothendieck, shelfmark 46, batch 2 (pages 21-40).
% Established from the facsimile archives/batches/46/batch-02.pdf (cut from a
% copy of 46.pdf fetched through the project relay,
% https://grothendieck-relay.onrender.com/source/46.pdf; PDF page k =
% archivists' page 20+k, no cover sheet), per the skill transcribe-grothendieck.
% The folder is ninety-four pages in five batches; this is the second.
%
% Pass: Opus 5.5 (claude-opus-5-5), 2026-09-26 — first pass, unchecked against the pages by a human.
%
% Pages skipped, and why:
%   - 22: a stencilled (blue) typescript, « CHAPITRE II. STRUCTURES
%     UNIFORMES », scanned upside down, the other face of the leaf of page 21;
%     unrelated to the mathematics in hand and carrying no ink of his.
%   - 34 and 36: two typed leaves (the second numbered « 11 ») on the
%     restriction maps H^n(X,F) -> H^n(Y,h(F)) for torsion sheaves, with
%     their arrows left blank for inking; they are the other faces of the
%     leaves of pages 33 and 35 (the typing shows through), i.e. scrap paper
%     for the Frobenius notes. No ink of his on them; skipped as unrelated
%     versos, per the brief. Their text was not identified against a
%     published source in this pass.
%   - 38: back of a folder (the tan sheet of page 37); its only mark is the
%     word « Steinberg », upside down at the foot; recorded in a \note with
%     page 37's cover.
% Covers, no \page marker, words as headings with a \note: 27 (« Frobeniuseries
% pour faisceaux étales »), 37 (inscription cancelled), 39 (« Brauer-Severi.
% Dans Chap VI »).
%
% Typescript (the folder's « tapuscrit (s.d.) »): pages 23-26, « Définition du
% morphisme Verschiebung », typed pages [1], -2, -3, -4. Decision: NOT his. The
% text is a letter-note addressed to him in the second person (« Je ne sais
% pas exactement ce que tu veux que je définisse. Cependant, comme tu affirmes
% avoir perdu mon tapis de l'année dernière, je vais essayer de te le
% revendre ») by a correspondent who is not named on these pages; it answers a
% request, which is exactly the request formulated on page 21 (margin: define
% the « Verschiebung » V_S(X)). The symbols (tensor signs, arrowheads, Delta)
% are inked in blue into blanks left by the typewriter, as the typist-author
% would; nothing on the four pages is identifiably an annotation of his. The
% author is not guessed. Each page therefore gets one French \note (content
% identified, no re-setting), per the brief.
%
% Pages summarised: none of his. The typescript pages 23-26 are described,
% not transcribed (see above). Page 28 (cancelled lines on a folder) is a
% \page with a single \note.
%
% His pagination: none on the manuscript leaves. The typescript pages carry
% « -2 », « -3 », « -4 » at top right.
%
% What the batch is about. Page 21 (continuing the Frobenius notes of batch 1,
% though not the local question on which page 20 breaks off): a problem sheet
% — define for X over S of char. p a Cartier-type operation on Omega^1_{X/S}
% transposed to D -> D^p on S-derivations, compatible with base change; in the
% margin, b) define a « Verschiebung » V_S(X). Pages 23-26: a colleague's
% typed answer defining V for a commutative affine k-group via the symmetric
% tensors Sigma^p A, with F.V = p and V.F = p. Pages 29-33, under the cover
% « Frobeniuseries pour faisceaux étales »: the relative Frobenius of an étale
% sheaf, Fr_{F/X} : Fr_X^*(F) -> F (and its adjoint Fr_{X*}(F) -> F), built
% from U -> U^{(p/X)} for U étale over X; Prop. 1) Fr_{F/X} is an isomorphism,
% functorial in X, F and compatible with base change; 2) characterisation of
% Frobenius-type maps by these conditions (for X = Spec F_p, Fr_X = id);
% « Théorème de trivialité cohomologique du Frobenius absolu »: Fr acts as the
% identity on H^0, H^1, and on RΓ_X of complexes in D(A); a formal proof via
% the fact that Fr_X^* on the étale topos is canonically isomorphic to the
% identity functor. Page 35 revisits the argument and notes its failure
% points. Page 40, under the cover « Brauer-Severi »: definition of a
% Brauer-Severi scheme P/S, Propositions 1-2 (smoothness, projectivity; the
% invertible sheaf omega^{-1}; P a projective bundle iff an invertible sheaf of
% relative degree 1 exists), a section implies projective bundle.
%
% Notation fixed once for the batch, as in batch 1: his double-struck F with
% r is set \mathrm{Fr} (Fr_X absolute, Fr_{X'/X} relative, Fr_{F/X} for a
% sheaf F); X^{(p/X)}, U^{(p/X)} as written; underlined sheaves
% (\underline{\Omega}, \underline{\omega}) keep their underline in math.
% The étale topos is X_{\text{ét}}, with a tilde where he writes one.
%
% Readings to check first: the whole of page 21 (fast, heavily worked); the
% margins of pages 21, 31 and 33; the cancelled top of page 35; page 32's
% « Corollaire 2 » line; page 40's margin.

\documentclass[11pt,a4paper]{article}
% Le préambule commun à toutes les transcriptions du fonds Grothendieck.
%
% Il tient en une idée : **l'incertitude se compose, elle ne se résout pas.**
% Une transcription qui lisse un mot douteux a détruit la seule chose pour
% laquelle on la faisait — un lecteur doit pouvoir savoir, à chaque endroit,
% ce qui était sur le feuillet et ce qui a été ajouté. Les six macros qui
% suivent sont donc l'appareil critique, et non de la décoration.
%
% Les mêmes commandes sont reconnues par `scripts/render.mjs`, qui produit la
% vue de lecture du site. Une macro ajoutée ici doit l'être aussi là-bas,
% faute de quoi le rendu échouera bruyamment — ce qui est voulu : un rendu
% silencieusement faux est pire qu'un rendu absent.

\ProvidesPackage{grothendieck}[2026/08/08 Transcriptions du fonds Grothendieck]

\usepackage{amsmath,amssymb,amsthm}
\usepackage{mathrsfs}
\usepackage{stmaryrd}
% Les diagrammes commutatifs sont partout dans ce fonds : c'est la langue
% même de la géométrie algébrique de Grothendieck, et une transcription qui
% les rendrait en prose ne transcrirait rien.
\usepackage{tikz-cd}
% Français par défaut — c'est la langue des feuillets — mais l'anglais est
% chargé aussi : la traduction et le résumé sont des documents anglais, et la
% typographie française (espaces avant les deux-points) y serait une faute.
% Ces documents basculent par \selectlanguage{english} en tête de corps.
\usepackage[english,french]{babel}
\usepackage{fontspec}
\usepackage{geometry}
\geometry{a4paper,margin=25mm,marginparwidth=28mm}
\usepackage{marginnote}
\usepackage{xcolor}
% ulem, not soul. `soul`'s \st and \ul are built on a character-by-character
% scan that XeLaTeX's font machinery breaks: the document fails with an
% undefined control sequence rather than degrading. ulem does the same two jobs
% and is Unicode-engine safe. `normalem` stops it hijacking \emph, which the
% transcriptions use for Grothendieck's own emphasis.
\usepackage[normalem]{ulem}
\usepackage{hyperref}
\hypersetup{colorlinks=true,linkcolor=black,urlcolor=black,pdfborder={0 0 0}}

% --- Métadonnées du lot -----------------------------------------------------
% Reprises telles quelles par le script de rendu, qui les affiche en tête de
% la vue de lecture. Elles fixent ce que le fichier prétend couvrir : sans
% elles, une transcription flotte sans point d'attache dans un fonds de seize
% mille feuillets.

\newcommand{\folder}[1]{\gdef\grothFolder{#1}}
\newcommand{\batch}[1]{\gdef\grothBatch{#1}}
\newcommand{\leaves}[2]{\gdef\grothFirst{#1}\gdef\grothLast{#2}}
\newcommand{\foldertitle}[1]{\gdef\grothFoldertitle{#1}}
\gdef\grothFolder{?}\gdef\grothBatch{?}\gdef\grothFirst{?}\gdef\grothLast{?}\gdef\grothFoldertitle{}

% --- Le résumé --------------------------------------------------------------
%
% Il ouvre la lecture modernisée, sous le titre « Résumé », et il est composé
% en petit. Ce n'est pas un ornement : c'est le seul passage du document qui
% ne soit pas des mathématiques, et un lecteur doit pouvoir le reconnaître
% avant de l'avoir lu — puis le sauter s'il connaît déjà le sujet, ou s'y
% arrêter s'il ne le connaît pas. Le filet qui le suit marque la reprise du
% corps.

\newenvironment{resume}
  {\small\setlength{\parskip}{0.45em}}
  {\par\medskip\hrule\medskip}

% --- Mots-clés ---------------------------------------------------------------
%
% En anglais, en clôture du résumé de la lecture modernisée : le vocabulaire
% moderne sous lequel chercher ce que les feuillets construisent. Le manifeste
% du site (`npm run manifest`) les extrait pour étiqueter la cote entière —
% cette ligne est la source unique des tags, il n'y a pas de fichier à côté.
\newcommand{\keywords}[1]{\par\smallskip\noindent
  {\footnotesize\textsc{Keywords} --- \textit{#1}}\par}

% --- L'appareil critique ----------------------------------------------------

% Le numéro de feuillet, en marge. C'est l'ancre qui permet de régler un
% désaccord sur une formule en regardant *un* feuillet plutôt qu'une chemise
% de six cents. Le site s'en sert aussi pour tourner les pages du fac-similé
% quand on fait défiler la transcription.
\newcommand{\leaf}[1]{%
  \marginnote{\footnotesize\color{gray!70}\textsf{#1}}%
  \label{leaf:#1}\ignorespaces}

% Illisible : on ne devine pas. Un mot inventé qui se lit comme les autres est
% le pire résultat possible.
\newcommand{\ill}{\textcolor{red!60!black}{[\,\ldots\,]}}

% Lecture douteuse : le mot est donné, et signalé comme incertain.
\newcommand{\uncertain}[1]{\uline{#1}}

% Ajout de l'éditeur — une lettre manquante, un mot sous-entendu.
\newcommand{\add}[1]{\textcolor{blue!50!black}{[#1]}}

% Biffé par Grothendieck lui-même. Ce qu'il a rayé fait partie du document :
% une rature dit souvent où la pensée a changé de direction.
\newcommand{\struck}[1]{\sout{#1}}

% Note du transcripteur, sur l'état matériel du feuillet.
\newcommand{\note}[1]{\par\smallskip{\small\color{gray!60!black}\textit{[#1]}}\par\smallskip}

% Note marginale de Grothendieck, distincte du corps du texte.
\newcommand{\marginal}[1]{\par\smallskip\noindent\hspace{1em}%
  {\small\color{gray!60!black}#1}\par\smallskip}

% --- Titre ------------------------------------------------------------------

\AtBeginDocument{%
  \begin{center}
    {\large\bfseries Fonds Alexandre Grothendieck --- cote n\textsuperscript{o}~\grothFolder}\\[2pt]
    {\itshape\grothFoldertitle}\\[4pt]
    {\small feuillets \grothFirst--\grothLast\ (lot \grothBatch)}\\[2pt]
    {\footnotesize\color{gray!60!black}Transcription établie d'après le fac-similé numérisé
     par l'Université de Montpellier.\\
     Les lectures incertaines sont soulignées, les illisibles notées [\,\ldots\,],
     les ajouts entre crochets.}
  \end{center}
  \vspace{1em}}

\endinput


\folder{46}
\batch{2}
\pages{21}{40}
\dating{[vers 1966-vers 1972]}
\watermark{Édition de démonstration}
\foldertitle{Schémas en groupes et fibrés principau[x] (Divers) : notes manuscrites (s.d.), tapuscrit (s.d.)}

\begin{document}

\page{21}

\emph{Pb} \quad $X$ sur $S$ de car.\ $p > 0$.

a) Définir une opération \uncertain{« de Cartier »}\ldots\ \uncertain{$p$-linéaire}
\[
  \underline{\Omega}^1_{X/S} \longrightarrow\ ?\ \struck{\ill{}}(\underline{\Omega}^1_{X/S})
\]
qui soit transposée dans un sens convenable \ill{} de
$\struck{\ill{}}\ \partial \mapsto \partial^{p}$ sur les $S$-dérivations, et
compatible avec les \uncertain{changements} de base

\begin{tikzcd}[column sep=small, row sep=small]
X \arrow[d] & X' \arrow[l] \arrow[d] \\
S & S' \arrow[l]
\end{tikzcd}

Il faudrait que lorsque \ill{} \ill{} \ill{} \ill{} formule, on trouve
\uncertain{quelque} chose dans \struck{\ill{}}
$\boxed{\underline{\Omega}^1_{X/S}{}^{(p)}}$ \uncertain{Symétrisons}
(\uncertain{Frobeniusienne} de $\underline{\Omega}^1_{X/S}$), ce qui
\uncertain{est} l'opération de Cartier. Si les $\omega_i$, $d_j$ sont duales, \ill{}
\ill{} \ill{} \ill{} par $d$, $[\,,\,]$, et \ill{} \ill{} relation
$\langle \omega_i, d_j \rangle = \ldots$ \ill{} soit : \ill{} \ill{}
\[
  \boxed{\mathrm{Symm}^{p}(\underline{\Omega}^1_{X/S})}\ .
\]

b) \uncertain{Montrer} \ill{} \ill{} \ill{} \ill{} \uncertain{infinitésimaux} \ill{}
\ill{} $1$ sur $S$ \ill{} \ill{} $\Omega$ \uncertain{qu.\ coh.}\ sur $S$, avec
\[
  C : \Omega \to \mathrm{Symm}^{p}_{\mathcal{O}_S}(\Omega), \qquad
  d_1 : \Omega \to {\textstyle\bigwedge}^{2}_{\mathcal{O}_S} \Omega
\]
\marginal{\ill{} \ill{} \uncertain{structures analogues} \ill{}}

\marginal{\ill{} \uncertain{identification} \ill{} $d_1 = 0$ \ldots\
« \uncertain{Verschiebung} » $V_S(X) : P_S(G) \to G$ \ldots\ \ill{} groupe
\uncertain{commutatif} $G$ sur $S$.}
\note{la page est écrite vite, à l'encre noire, sur le verso d'une feuille
ronéotypée ; la moitié des mots de liaison ne se lit pas. La note
« Verschiebung » est écrite verticalement dans la marge gauche, en trois lignes
très effacées, avec un « b) » ; la notation $P_S$ est celle de la page 11. La
demande de définir la Verschiebung est celle à laquelle répond le tapuscrit
des pages 23 à 26}

\page{23}

\note{Tapuscrit d'une autre main, « Définition du morphisme Verschiebung »,
adressé à Grothendieck au tutoiement (« Je ne sais pas exactement ce que tu
veux que je définisse. Cependant, comme tu affirmes avoir perdu mon tapis de
l'année dernière, je vais essayer de te le revendre ») ; l'auteur n'est pas
nommé. Page 1 : base le spectre d'un corps $k$ de caractéristique $p > 0$ ; cas
d'un $k$-groupe commutatif affine $G$ d'anneau $A$ ; notations $A^{(p)}$,
$F : A^{(p)} \to A$, $G^{(p)}$, $\otimes^{p}A$, $\Delta^{p}$, $m^{p}$. Les
symboles ($\otimes$, $\Delta$, pointes de flèches) sont ajoutés à l'encre bleue
dans les blancs laissés par la machine ; aucune annotation de sa main}

\page{24}

\note{Tapuscrit, suite (p.~-2) : $\Sigma^{p}A$ (tenseurs symétriques),
$S^{p}A$ (image dans l'algèbre symétrique), les applications
$i : \Sigma^{p}A \to \otimes^{p}A$, $q : \otimes^{p}A \to S^{p}A$,
$b : S^{p}A \to A$, $r : \Sigma^{p}A \to A^{(p)}$ et
$j : A^{(p)} \to S^{p}A$, calculées sur une base de $A$. Aucune annotation de
sa main}

\page{25}

\note{Tapuscrit, suite (p.~-3) : $q.i = j.r$ (renvoi à Serre, \emph{Groupes
algébriques et corps de classes}, preuve du lemme 11, p.~62) ; diagramme
commutatif d'algèbres ; l'homomorphisme $\underline{r}$ ; définition de la
Verschiebung $V$ par $F.V = \underline{p}$ ; $F$ et $V$ s'échangent dans la
dualité de Cartier ; passage aux variétés algébriques. Aucune annotation de sa
main}

\page{26}

\note{Tapuscrit, fin (p.~-4) : diagramme
$G \to \Sigma^{p}G \to G^{(p)}$ (diagonale, symétrisation) ; « Le diagramme
ci-dessus est donc défini pour tout groupe algébrique commutatif, affine ou
non. On posera donc $V = a.r$ et on aura $V.F = p$. » Aucune annotation de sa
main}

\section*{Frobeniuseries pour faisceaux étales}

\note{inscrit et souligné en haut à droite d'une chemise (page 27), qui ne
reçoit pas de numéro de page}

\page{28}

\note{au bas du feuillet (le dos de la chemise), quatre lignes de sa main
barrées d'un quadrillage de traits obliques, avec « XV » \uncertain{} souligné
à droite ; illisibles sous les ratures}

\page{29}

\emph{\uncertain{Proposition}}

Considérons un préschéma $X/\mathbb{F}_p$, et
\[
  X \xleftarrow{\ \mathrm{Fr}_X\ } X
\]
Comme $\mathrm{Fr}_X$ est un \uncertain{homéom.\ univ.}, on conclut que pour un
faisceau $F$ sur $X$ et un faisceau $G$ sur $X$ (\uncertain{pour top.\ étale}),
la donnée de $G \xrightarrow{\ \alpha\ } \mathrm{Fr}_X^{*}(F)$
\uncertain{équivaut} à la donnée de $\mathrm{Fr}_{X*}(G) \xrightarrow{\ \beta\ }
F$ (car $\mathrm{Fr}_{X*}$ et $\mathrm{Fr}_X^{*}$ sont \uncertain{inverses}
\ill{} foncteurs adjoints \ill{} \ill{} \uncertain{habituel}).

Ainsi, la donnée de $\alpha$ \ill{}
\[
  \mathrm{Fr}_{X*}(G) \xrightarrow{\ \mathrm{Fr}_{X*}(\alpha)\ }
  \mathrm{Fr}_{X*}\mathrm{Fr}_X^{*}(F) \xrightarrow[\ \sim\ ]{\ \mathrm{can}^{-1}\ } F ,
\]
et \uncertain{inversement} la donnée de $\beta$ définit
\[
  G \xrightarrow[\ \sim\ ]{\ \mathrm{can}^{-1}\ }
  \mathrm{Fr}_X^{*}\mathrm{Fr}_{X*}(G) \xrightarrow{\ \mathrm{Fr}_X^{*}(\beta)\ }
  \mathrm{Fr}_X^{*}(F) .
\]
Nous \uncertain{allons} \struck{\uncertain{identifier}} \ill{} \ill{} \ill{} \ill{}
\ill{} \ill{} \ill{} \ill{} \ill{} \ill{} \ill{} \ill{}
\note{deux lignes entourées, d'une écriture très petite, ne se lisent pas}
\ill{} \ill{} \ill{} \ill{} $= G$, \uncertain{définir} un hom.\
\emph{\uncertain{canonique}}
\begin{align*}
  &(1) \qquad F \xleftarrow{\ \mathrm{Fr}_{F/X}\ } \mathrm{Fr}_X^{*}(F) \qquad \text{i.e.} \\
  &(2) \qquad \mathrm{Fr}_{X*}(F) \xrightarrow{\ \mathrm{Fr}_{F/X*}\ } F
\end{align*}
(qui \uncertain{sera} d'ailleurs un isom., mais \ill{} \ill{} \ill{}
\ill{} \ill{} \ill{} cas !). \ill{} \ill{} \ill{} (2) \ill{} \ill{}
\uncertain{définition} \uncertain{faisceau}\ldots
\note{une accolade dans la marge gauche embrasse la première moitié de la
page ; la flèche de (1) a une pointe à chaque bout, la pointe droite plus
appuyée}

\page{30}

\ill{} \ill{} \ill{} \uncertain{évident}. \ill{} \ill{}, si $U$ \uncertain{est}
étale sur $X$
\[
  \mathrm{Fr}_{X*}(F)(U) = F(\mathrm{Fr}_X^{*}(U)) = F(U^{(p/X)})
\]
\struck{or $\mathrm{Fr}_X^{*}(U) \simeq U^{(p/X)}$} or \ill{} \ill{} \ill{}
\emph{isom} canonique fonctoriel en $U$
\[
  U \xrightarrow[\ \sim\ ]{\ \mathrm{Fr}_{U/X}\ } U^{(p/X)}
\]
qui donne un isom.\ fonctoriel en $U$
\[
  \mathrm{Fr}_{X*}(F)(U) = F(U^{(p/X)}) \xrightarrow{\ \sim\ } F(U)
\]
\ill{} qui \uncertain{est} l'hom.\ cherché
\[
  \boxed{\mathrm{Fr}_{F/X*} : \mathrm{Fr}_{X*}(F) \longrightarrow F}
\]
\note{dans la marge gauche, un petit carré : $U \to X$ au-dessus de
$U \to X$, flèches verticales $\mathrm{Fr}_X$ et $\mathrm{Fr}_U$}

\struck{Supposons que $F$ soit représenté par $X'$ étale sur $X$. \ill{}
\ill{} \ill{} \ill{} \ill{}}
\struck{Considérons \ill{} \ill{} correspondante}
\[
  \struck{\mathrm{Fr}_{F/X}^{*} : \ill{}\,F \to \mathrm{Fr}_X^{*}F} \qquad
  \struck{\mathrm{Fr}_X^{*}\mathrm{Fr}_{X*}(F) \nearrow}
\]
\struck{Si $U$ est un \ill{} étale sur $X$, \ill{} définir
$F(U) \to \mathrm{Fr}_X^{*}(F)(U)$ \ill{} \ill{} $U$ \ill{} \ill{} de
$\mathrm{Fr}_X^{*}(V)$, $V$ étale \ill{} $X$, \ill{} \ill{}
$\mathrm{Fr}_X^{*}\mathrm{Fr}_{X*}(F)(U) \xrightarrow{\sim}
\mathrm{Fr}_{X*}(F)(V)$}
\note{ce passage, dans un cadre, est barré de grands traits obliques}

\emph{Prop} 1) $\mathrm{Fr}_{F/X}^{*} : \mathrm{Fr}_X^{*}(F) \longrightarrow F$
\uncertain{pour} $X$, $F$ variables \uncertain{est} compatible aux
\uncertain{extensions} de la base quelconques \ill{}
\uncertain{commutativité} dans

\begin{tikzcd}[column sep=small, row sep=small]
X \arrow[d, "\mathrm{Fr}_X"'] & X' \arrow[l, "g"'] \arrow[d, "\mathrm{Fr}_{X'}"] \\
X & X' \arrow[l, "g"]
\end{tikzcd}

$X' \xrightarrow{\ g\ } X$, i.e.\ \ill{}
\[
  g^{*}(\mathrm{Fr}_X^{*}(F)) \xleftarrow{\ g^{*}(\mathrm{Fr}_{F/X}^{*})\ }
  g^{*}(F) \overset{\text{df}}{=} F'
\]
\[
  \mathrm{Fr}_{X'}^{*}(g^{*}(F)) = \mathrm{Fr}_{X'}^{*}(F')
  \xleftarrow{\ \mathrm{Fr}_{F'/X'}^{*}\ } F'
\]
\note{les deux lignes sont reliées par « $\simeq$ trans.\ » entre
$g^{*}(\mathrm{Fr}_X^{*}(F))$ et $\mathrm{Fr}_{X'}^{*}(g^{*}(F))$, et une flèche
oblique va de $F'$ vers la première ligne ; le sens des flèches est tel
qu'écrit, en désaccord avec celui de l'énoncé}

\page{31}

D'autre part, si $F$ est un faisceau constant, de la forme $I_X$, \uncertain{on a}
\[
  \mathrm{Fr}_X^{*}(I_X) \simeq I_X
  \xleftrightarrow{\ \mathrm{Fr}_{F/X}^{*} = \mathrm{id}_{I_X}\ } I_X
\]

2) \struck{Ces conditions} \struck{\ill{} $X$ \ill{} il suffit \ill{}
\ill{}} Ces conditions (pour $X$, $F$ variables) \ill{} $X'$ étale sur $X$ \ill{}
et la fonctorialité en $F$, \emph{caractérisent} les hom.\
$\mathrm{Fr}_{F/X}$.

\marginal{NB. La \uncertain{condition} \ill{} : $I_X$ \ill{} \ill{} \ill{}
\uncertain{identité}, on \ill{} \ill{} \ill{} $I$ \ill{} \ill{} \ill{}
\ill{} \ill{} \ill{} \ill{}, on \ill{} \ill{} faisceau \ill{} \ill{}
\uncertain{endomorphismes} !}\note{dans la marge gauche, de biais}

Prouvons l'assertion 2°). Comme \uncertain{tout} $F$ \ill{} \ill{}
\uncertain{quotient} d'un $G$ représentable, \ill{} \ill{} \ill{} \ill{}

\begin{tikzcd}[column sep=small, row sep=small]
\mathrm{Fr}_X^{*}(G) \arrow[d] & G \arrow[l] \arrow[d, "\text{épim}"] \\
\mathrm{Fr}_X^{*}(F) & F \arrow[l]
\end{tikzcd}

\ill{} \uncertain{déterminer} $\mathrm{Fr}_{F/X}^{*}$ pour $F$ représentable.
La condition de fonctorialité \ill{} $X$ \struck{\uncertain{variable}} (\ill{}
\ill{} $F$) \uncertain{permet} de nous \uncertain{localiser} sur $F$ et sur $X$,
ce qui nous ramène au cas où $F \to X$ est une \uncertain{immersion ouverte},
\uncertain{puis} un isom., qui est justiciable de la $2^{\text{e}}$ condition
énoncée dans 1).

\marginal{$2$}\note{un grand « 2 » est écrit dans la marge gauche, en regard
du paragraphe suivant}

\uncertain{Supposons} $\mathrm{Fr}_X = \mathrm{id}$ (ce qui signifie que $X$ est
défini par un \uncertain{anneau} loc.\ \ill{} \ill{} discontinu --- \ill{}
\ill{} \ill{} \ill{} \ill{} \ill{} \ill{} faisceau \ill{} \ill{}
\ill{} : \ill{} \ill{} $\mathbb{F}_p$ ---) p.\ ex.\
$X = \operatorname{Spec} \mathbb{F}_p$. Alors \ill{} \ill{} \ill{} \ill{}
\ill{} \ill{} \ill{}
$\mathrm{Fr}_X^{*}(F) = \mathrm{id}_X^{*}(F) \simeq F$, et
$\mathrm{Fr}_{F/X}$ \uncertain{devient} \ill{} \emph{\uncertain{automorphisme}} de
$F$. \ill{} \ill{} \ill{} \ill{} \uncertain{évident} : \ill{} \ill{}
l'\uncertain{identité} \ill{} \ill{} $F$ \ill{} \ill{} sur
$X = \operatorname{Spec} \mathbb{F}_p$\ldots
\note{la fin de la page, serrée au bas du feuillet, est en partie
surchargée}

\page{32}

\emph{Prop} \uncertain{Autre caractérisation} de $\mathrm{Fr}_{F/X}^{*}$, pour
$X$ fixé :
\begin{quote}
c'est un hom.\ fonctoriel en $F$, et pour $F$ \emph{représentable}
\add{\uncertain{par $X'$ étale$/X$}}, c'est l'hom.\ de Frobenius habituel noté
$\mathrm{Fr}_{X'/X}$.
\end{quote}

\emph{Théorème de trivialité cohomologique du Frobenius absolu.}
\begin{quote}
L'hom.
\[
  H^0(X, F) \longrightarrow H^0(X, F), \qquad
  H^0(X, F) \simeq H^0(X, \mathrm{Fr}_{X*}(F))
  \xrightarrow{\ H^0(\mathrm{Fr}_{F/X})\ } H^0(X, F)
\]
est l'\emph{identité}, pour tout faisceau $F$ sur $X_{\text{ét}}$.
\end{quote}

\emph{Corollaire 1} Si $F$ est un faisceau en groupes,
\[
  H^1(X, F) \simeq H^1(X, \mathrm{Fr}_{X*}(F))
  \xrightarrow{\ H^1(\mathrm{Fr}_{F/X})\ } H^1(X, F)
\]
est l'\emph{identité}.

\emph{Corollaire 2} Si $F^{\bullet}$ est un
\add{\uncertain{complexe de}} faisceaux \struck{\uncertain{abéliens}}
\add{\uncertain{de $A$-Modules}}, (\struck{$A$ un faisceau \ill{} \ill{}
\ill{}} $A$ un \uncertain{anneau} \ill{} \ill{}),
\[
  R\Gamma_X(F^{\bullet}) \simeq R\Gamma_X(\mathrm{Fr}_{X*}(F^{\bullet}))
  \xrightarrow{\ R\Gamma_X(\mathrm{Fr}_{F^{\bullet}/X})\ } R\Gamma_X(F^{\bullet})
\]
est l'\emph{identité} \add{\uncertain{dans}} la catégorie dérivée $D(A)$ de la
catégorie des $A$-modules. En particulier $H^i(X, F) \to H^i(X, F)$ est
l'identité.
\note{dans chacun des trois énoncés, la flèche oblique
$H^0(\mathrm{Fr}_{F/X})$ (resp.\ $H^1$, $R\Gamma_X$) part de la ligne inférieure
et aboutit au second terme de la ligne supérieure ; on l'a mise en ligne}

Les corollaires résultent du th.\ en utilisant le calcul de $R^1\Gamma_X$
\uncertain{resp.}\ de $R\Gamma_X$ par \uncertain{résolutions}. / Pour prouver le
th., on \ill{} \struck{\ill{} \ill{} \ill{} \ill{} \ill{} \ill{}}
\struck{\ill{} \ill{} \ill{} \ill{}} l'hom.
\[
  H^0(X, F) \xrightarrow{\ H^0(\mathrm{Fr}_{F/X})\ } H^0(X, F)
\]
\ill{} \emph{fonctoriel} en $F$. \uncertain{Quand} \ill{}, \ill{} \ill{}
\ill{}, c'est l'identité, \uncertain{puisque} si $F$ \ill{} le faisceau
\ill{} \ill{} \ill{} \ill{} \ill{} \uncertain{trivial}\ldots\ \ill{} \ill{}
\ill{} \ill{} \ill{} \ill{} \ill{} \uncertain{condition}.

\page{33}

$\mathbb{S}$ \uncertain{catégorie} \ill{} \uncertain{produits fibrés}.
$\mathcal{M}$ \uncertain{ens.}\ de morphismes \ill{} \ill{} \ill{} \ill{}, \ill{}
\uncertain{changements} de base (\uncertain{un} \ill{} \ill{}), défini pour tout
$X \in \mathrm{Ob}\,\mathbb{S}$ une catégorie $X_{\text{ét}}$. (\ill{} \ill{}
\uncertain{topologie} \uncertain{standard}),

\uncertain{Supposons} \uncertain{que} $X_{\text{ét}}$, \ill{} \ill{} \ill{} \ill{}
\ill{} \ill{} \ill{} pour tout $f : X \to Y$ dans $\mathbb{S}$, \ill{}
foncteur
\[
  Y_{\text{ét}} \xrightarrow{\ f^{*}\ } X_{\text{ét}}
\]
\uncertain{compatible} \ill{} « \uncertain{cartésien} ».
\uncertain{Considérons} \struck{\ill{}} \uncertain{un} \ill{} $\mathbf{F}$ endom.\
du foncteur identique de $\mathbb{S}$.

a) \uncertain{Pour} tout $X \in \mathbb{S}$, $\mathbf{F}_X : X \to X$ \ill{}
\uncertain{une} \uncertain{équivalence} \ill{}
$\mathbf{F}_X^{*} : \widetilde{X}_{\text{ét}} \to \widetilde{X}_{\text{ét}}$ \ill{}
\ill{} \uncertain{et} $p : X' \to X$ \uncertain{dans} $\mathcal{M}$, \ill{} \ill{}
$X' \to X'^{(p/S)}$ \ill{} \ill{} \ill{}, i.e.\ le \ill{} :

\begin{tikzcd}[column sep=small, row sep=small]
X' \arrow[d, "p"'] & X'^{(p/X)} \arrow[l] \arrow[d] & X' \arrow[l, "\mathrm{Fr}_{X'/X}"'] \arrow[dl, "p"] \arrow[ll, bend right=30, "\mathbf{F}_{X'}"'] \\
X & X \arrow[l, "\mathbf{F}_X"] &
\end{tikzcd}

\begin{tikzcd}[column sep=small, row sep=small]
X' \arrow[d, "\mathbf{F}_{X'}"'] & X' \arrow[l, "\mathbf{F}_{X'}"'] \arrow[d, "p"] \\
X & X \arrow[l, "\mathbf{F}_X"]
\end{tikzcd}

\uncertain{est} \uncertain{cartésien}.

\marginal{\uncertain{b) $\Longleftrightarrow$ que}}\note{entouré, au-dessus de
« Dém. »}
\emph{Dém.} \uncertain{Le} foncteur
$\mathbf{F}_X^{*} : X_{\text{ét}} \to X_{\text{ét}}$ \uncertain{est isomorphe}
\uncertain{au} \uncertain{foncteur} \uncertain{identique}, \uncertain{d'où}
\uncertain{résulte} que le \uncertain{morphisme}
$\mathbf{F}_X^{*} : \widetilde{X}_{\text{ét}} \to \widetilde{X}_{\text{ét}}$
\uncertain{est une} \uncertain{équivalence}, \uncertain{qui} \uncertain{induit}
\ill{} \uncertain{isomorphisme} \ill{} \ill{} \ill{} \uncertain{foncteurs}
\uncertain{identiques} \ill{}, \ill{} \ill{} \ill{} \uncertain{isomorphismes}
\uncertain{donnés} \uncertain{par} \ill{} \ill{} $\mathbf{F}_{-/X}$ sur les
\uncertain{faisceaux} \uncertain{représentables}. Donc b) $\Longrightarrow$ a).
\marginal{\ill{} \uncertain{de Topos}}\note{à droite de la ligne de
« Dém. », souligné, avec un signe de vérification}
\note{la page est de la même encre bleue que les précédentes, sur le verso
d'un feuillet dactylographié. La lettre notée ici $\mathbf{F}$ (un F double,
proche de son $\mathrm{Fr}$ sans le r) désigne un endomorphisme du foncteur
identique qu'il distingue du Frobenius ; lecture de la notation incertaine}

\page{35}

\struck{$\mathbf{F}_{\alpha}^{*}(\mathbf{F}_{X}^{*}(X')) \longrightarrow
\mathbf{F}_{\alpha}^{*}(X')$}\quad
\struck{$\mathbf{F}_{X}^{*}(F)$}
\note{deux lignes en haut à gauche, raturées ; au-dessus de la première,
« $\mathbf{F}_{\alpha'}$ » et une flèche courbe}

\marginal{$\mathbf{F}_X$ ?}\note{un signe dans la marge, à gauche du
paragraphe suivant, peut-être un renvoi}
\uncertain{Peut-être} que a) \uncertain{implique} aussi b), \ill{} \ill{} \ill{}
\uncertain{s'en} \ill{} \ill{} \ill{} \uncertain{renforce} a) \ill{}
\uncertain{exigeant} que $\mathbf{F}_X$ \uncertain{commute} \ill{}
\uncertain{propriétés} dites \ill{} \uncertain{tout} changement de base (\ill{}
\ill{} \ill{}, \ill{} \ill{} \ill{} \ill{} $\mathbf{F}_X : X \to X$ !).

\begin{tikzcd}[column sep=small, row sep=small]
X' \arrow[d, "p"'] & X'^{(p/S)} \arrow[l] \arrow[d] & X' \arrow[l, "\sim"'] \arrow[dl, "p"] \\
X & X \arrow[l, "\mathbf{F}_X"] &
\end{tikzcd}

\note{au-dessus du diagramme, « $\exists$ isom » et une flèche courbe
$\mathbf{F}_{X'}$, barrée}

La \uncertain{définition} de $\mathbf{F}_{X'/X}^{*}$
\add{\uncertain{$(X' \in \mathrm{Ob}\,X_{\text{ét}})$}} (\uncertain{comme}
\ill{}) \uncertain{faisceaux} \ill{} $\mathbf{F}_{F/X}$ \ill{} \ill{} \ill{}
\uncertain{fonctoriel} en $F$ \ill{} \ill{} \ill{}
\uncertain{conditions}, \ill{} \ill{} \ill{} \ill{} \ill{} \ill{}
\uncertain{identité} \ill{} \ill{} \uncertain{propriétés}).
\ill{} \ill{} \ill{} \uncertain{faisceaux} \uncertain{constants} :
\[
  \mathbf{F}_X^{*}(I_{X'}) = I_X \xrightarrow[\ \mathrm{id}_{I_X}\ ]{\ \mathbf{F}_{F/X}^{*}\ } I_X
\]
\ill{} \ill{} \ill{} \uncertain{en} \uncertain{remplaçant} \ill{} \ill{} \ill{}
\ill{} \ill{}. \uncertain{Car} \ill{} \struck{\uncertain{fonctorialité}}
$I \to$ \uncertain{identité}\ldots\ \uncertain{compatibilité} \ill{}
$\mathbf{F}_{F/X}$ par changement de base \uncertain{compatibilité} \ill{}
changement de base \ill{} \emph{\uncertain{quelconques}}. Le
\uncertain{compatibilité} \ill{} \ill{} \ill{} $F$, \ill{} : la
\uncertain{caractérisation}.
\begin{quote}
\ill{}, + \uncertain{fonctorialité} \ill{} \ill{} \ill{}

\uncertain{également} \struck{\ill{} \ill{} \ill{} \ill{} \ill{} \ill{}}
\end{quote}
\struck{\ill{} \ill{} \ill{} \ill{} \ill{}}
\note{les deux lignes réunies par un trait vertical à gauche sont en partie
biffées d'un long trait ondulé}
Il \uncertain{implique} l'identité sur les $H^0$, \ill{} \ill{} \ill{} $H^i$,
$R\Gamma_X$\ldots

\note{chemise (page 37) : en haut à droite, un titre souligné, barré et
illisible, puis trois lignes, « Variances \ill{} / \ill{} $G$-\ill{}-\ill{} /
(\ill{} \ill{}) », barrées d'un quadrillage de traits obliques ; le feuillet ne
reçoit pas de numéro de page. Au dos (page 38), seul, à l'envers au bas du
feuillet : « Steinberg »}

\section*{Brauer-Severi}

\note{inscrit et souligné en haut d'une chemise (page 39), avec au-dessous
« Dans Chap \emph{VI} » ; le feuillet ne reçoit pas de numéro de page}

\page{40}

\struck{\ill{}}\quad \emph{\struck{Systèmes linéaires et \ill{}}}
\emph{\add{Schémas} de Brauer-Severi}

\emph{Déf} \emph{Schéma de Brauer-Severi} \struck{\ill{}} $P/S$ : il devient
un fibré projectif loc.\ \uncertain{fppf}.

\struck{Remarques}

\emph{Proposition 1} Un tel $P$ est \uncertain{projectif} et lisse sur $S$.
Le faisceau \struck{\ill{}} $\underline{\omega}^{-1}_{X/S}$ est ample rel.\ à
$S$.

\emph{Proposition 2} \struck{Si $P/S$ \ill{} \ill{} \ill{} une section, alors
$P$} Pour que $P$ soit un fibré projectif, il \struck{faut et suffit} faut et
suffit qu'il existe un $\underline{P}$ \ill{} \ill{} de degré $1$ sur
\ill{} faisceau inversible \struck{\ill{} \ill{}} de degré $1$ sur chaque fibre,
\ill{} \ill{} \ill{} \ill{} \ill{} $=$, \ill{} \ill{} \ill{}
$\underline{L}^{\otimes(r+1)} \otimes \underline{\Omega}^{r}_{X/S}$ soit
\uncertain{trivial} : \ill{} \ill{} \uncertain{l'image} \uncertain{inverse}
d'un faisceau \uncertain{inv.}\ sur $S$.

\quad Cela signifie \ill{} que le \uncertain{résultat} de classification donné
au \ill{} I pour $\mathrm{Pic}(P)$ \struck{\ill{}} \uncertain{reste}
valable\ldots\ \emph{Corollaire} \uncertain{Les} \ill{} \ill{} \ill{}
\ill{}\ \uncertain{signification}.

\marginal{Construction de $\underline{\mathrm{Div}}^{1}_{P/S} = \check{P}$ et
\ill{} \ill{} \uncertain{projective} / \emph{Déf} $P/S$ \uncertain{projectivement}
\ill{}, \ill{} \ill{} \ill{} \ill{} \ill{}}\note{dans la marge gauche,
écrit verticalement, dans deux cadres reliés par une accolade au texte des
propositions}

[\emph{Proposition}] Si $P$ admet une section \add{sur $S$}, c'est un fibré
projectif : utiliser le théorème de descente \add{(et dualité des
projectifs)}.

\ldots\ \uncertain{Toutefois} \ill{} \ill{} si \uncertain{tout} Module loc.\
libre sur $S$ est libre (p.\ ex.\ $S$ local) \struck{\ill{} \ill{} \ill{}
\ill{}} \struck{\ill{} \ill{} \ill{}}.
\note{la dernière ligne est barrée en deux reprises, avec au-dessus
« \uncertain{principal} » ; la page s'arrête ici}

\end{document}
